
\paragraph{Law of large numbers.} We proved the weak form in \cref{thm:wlln-0908}.

\begin{theorem}[Law of Large Numbers]\label{thm:slln-0910}
    Let $X_1, X_2, \dots$ be independent, identically distributed (i.i.d.)\ random variables
    distributed as $X$, where $\expectation{X}$ is finite. Then, with probability $1$,
    \[\lim_{n \to \infty} \frac{\sum_{i=1}^{n} X_i}{n} = \expectation{X}.\]
\end{theorem}

\section{Random Walks on Graphs}

Let $G = (V, E)$ be a finite undirected graph without isolated vertices, with
$V = \set{1, \dots, n}$ and $m = \abs{E}$ edges. We write $u \sim v$ if $\set{u,v} \in E$, and
$\deg(u)$ (or $d_u$) for the number of neighbours of $u$. The walker is now called the token, and
$X_t$ is its vertex at time $t = 0, 1, 2, \dots$.

\begin{definition}[Distribution at Time $t$, Transition Matrix]\label{def:walk-0910}
    The \textbf{distribution at time $t$} is the row vector
    \[p_t = \left(p_t^1, p_t^2, \dots, p_t^n\right), \qquad p_t^i = \prob{X_t = i}.\]
    The \textbf{transition matrix} is the $n \times n$ matrix $P$ with entries
    \[P_{uv} = \prob{u \to v} = \prob{X_{t+1} = v \given X_t = u}
        = \begin{cases}
            \frac{1}{\deg(u)} & \text{if } u \sim v, \\
            0                 & \text{otherwise.}
        \end{cases}\]
\end{definition}

\subsection{The Triangle}

In the triangle on the vertices $1, 2, 3$ (the right half of \cref{fig:walks-0908}) every vertex has
degree $2$, so $P$ is the $3 \times 3$ matrix with $0$ on the diagonal and $\sfrac{1}{2}$ in every
other entry.

\subsection{The Evolution \texorpdfstring{$p_{t+1} = p_t P$}{p(t+1) = p(t) P}}

\begin{lemma}[Evolution of the Distribution]\label{lem:evolution-0910}
    For every $t \geq 0$,
    \[p_{t+1} = p_t P, \qquad \text{and so} \qquad p_t = p_0 P^t.\]
\end{lemma}

\begin{proof}
    Fix a vertex $v$. The events $\set{X_t = u}$, $u \in V$, are disjoint and one of them occurs, so
    \[p_{t+1}^v = \prob{X_{t+1} = v} = \sum_{u \in V} \prob{X_t = u \text{ and } X_{t+1} = v}
        = \sum_{u \in V} p_t^u P_{uv} = (p_t P)_v.\]
    Here $\prob{X_t = u \text{ and } X_{t+1} = v} = p_t^u P_{uv}$ by the definition of $P_{uv}$ as a
    conditional probability when $p_t^u > 0$, and both sides are $0$ when $p_t^u = 0$. So
    $p_{t+1} = p_t P$. Induction on $t$ (with $P^0 = I$) gives $p_t = p_0 P^t$.
\end{proof}

Entry by entry, $p_{t+1} = p_t P$ is \cref{eq:walkstep-0908}. For the triangle with $p_0 = (1,0,0)$,
$p_1$ and $p_2$ are in \cref{exa:walks-0908}.

\subsection{Row Stochastic Matrices}

\begin{definition}[Probability Vector, Row Stochastic Matrix]\label{def:stochastic-0910}
    A vector $p = (p^1, \dots, p^n)$ is a \textbf{probability vector} if $p^u \geq 0$ for all $u$ and
    $\sum_{u=1}^{n} p^u = 1$. An $n \times n$ matrix $P$ is \textbf{row stochastic} if every row of $P$ is a
    probability vector:
    \[P_{uv} \geq 0 \text{ for all } u, v, \qquad \text{and} \qquad \sum_{v=1}^{n} P_{uv} = 1 \text{ for all } u.\]
\end{definition}

The transition matrix of a random walk on any graph is row stochastic: row $u$ has $\deg(u)$ entries
equal to $\sfrac{1}{\deg(u)}$, and its other entries are $0$.

\needspace{9\baselineskip} % keep the exercise and the start of its solution on one page
\begin{exercisebox}
\begin{exercise}\label{ex:probvec-0910}\label{thm:probvec-0910}
    Let $P$ be a row stochastic $n \times n$ matrix and let $p$ be a probability vector. Show that
    $pP$ is a probability vector. So if $p_0$ is a probability vector and $p_{t+1} = p_t P$ for all
    $t \geq 0$, then $p_t = p_0 P^t$ is a probability vector for every $t \geq 0$.
\end{exercise}
\tcblower
\begin{proof}
    Let $q = pP$, so $q^v = \sum_{u=1}^{n} p^u P_{uv}$ for every $v$. Every term is a product of
    two nonnegative numbers, so $q^v \geq 0$. Every row of $P$ and the entries of $p$ sum to $1$, so
    $\sum_{v} q^v = \sum_{v} \sum_{u} p^u P_{uv} = \sum_{u} p^u \left(\sum_{v} P_{uv}\right)
        = \sum_{u} p^u = 1$.
    So $pP$ is a probability vector. The vector $p_0$ is one by assumption, and if $p_t$ is one,
    then so is $p_t P = p_{t+1}$. By induction on $t$, every $p_t$ is a probability vector.
\end{proof}
\end{exercisebox}

\subsection{The Transition Matrix of a General Graph}

Let $A$ be the \emph{adjacency matrix} of $G$, with $A_{uv} = 1$ if $u \sim v$ and $A_{uv} = 0$
otherwise. Let $D$ be the diagonal matrix with diagonal entries $d_1, \dots, d_n$. Its inverse
$D^{-1}$ has diagonal entries $\sfrac{1}{d_1}, \dots, \sfrac{1}{d_n}$ and exists because there are
no isolated vertices. Then $(D^{-1} A)_{uv} = \sfrac{A_{uv}}{\deg(u)}$, which is $P_{uv}$ of
\cref{def:walk-0910}, so $P = D^{-1} A$.
If $G$ is \emph{$d$-regular}, meaning every vertex has the same degree $d$, then $D = d I$ and
$P = \sfrac{A}{d}$. The triangle is $2$-regular, and its $P$ is $\sfrac{A}{2}$.

\section{Stationary Distribution}

\begin{definition}[Stationary Distribution]\label{def:stationary-0910}
    A probability vector $\pi$ is a \textbf{stationary distribution} of the walk with transition matrix $P$
    if
    \[\pi P = \pi \cdot 1 = \pi.\]
    Then $\pi$ is a left eigenvector of $P$ with eigenvalue $1$.
\end{definition}

By \cref{lem:evolution-0910}, if $p_t = \pi$ then $p_{t+1} = \pi P = \pi$. Only the neighbours $v$
of $u$ have $P_{vu} \neq 0$, and for them $P_{vu} = \sfrac{1}{\deg(v)}$, so (as in
\cref{eq:walkstep-0908})
\begin{equation}\label{eq:piP-0910}
    (\pi P)_u = \sum_{v \in V} \pi_v P_{vu} = \sum_{v \sim u} \frac{\pi_v}{\deg(v)}.
\end{equation}
The sum runs over the neighbours of $u$, as in \cref{fig:star-0910}.

\begin{figure}[ht]
    \centering
    \begin{tikzpicture}[scale=0.82]
        \node[point, label=left:{\small $u$}] (u) at (0,0) {};
        \node[point, label=above:{\small $v$}] (v) at (1.3,1.1) {};
        \node[point] (w1) at (-1.2,-1.2) {};
        \node[point] (w2) at (1.2,-1.2) {};
        \draw (u) -- (v);
        \draw (u) -- (w1);
        \draw (u) -- (w2);
    \end{tikzpicture}
    \caption{A vertex $u$ and its neighbours; $v$ is one of them.}
    \label{fig:star-0910}
\end{figure}

\begin{theorem}[Stationary Distribution of a Random Walk on a Graph]\label{thm:stationary-0910}
    Let $G$ be a graph with $m$ edges and no isolated vertices. Then
    \[\pi_v = \frac{\deg(v)}{2m} \qquad (v \in V)\]
    is a stationary distribution of the random walk on $G$.
\end{theorem}

\begin{proof}
    \emph{$\pi$ is a probability vector.} There are no isolated vertices and $V$ is not empty, so
    $m \geq 1$. So $\pi$ is well defined, and $\pi_v \geq 0$. Every edge $\set{u,v}$ is counted
    once in $\deg(u)$ and once in $\deg(v)$. So $\sum_{v \in V} \deg(v) = 2m$ (the handshake lemma)
    and $\sum_{v \in V} \pi_v = 1$.

    \emph{$\pi P = \pi$.} For every vertex $u$, by \cref{eq:piP-0910} and $\deg(v) \geq 1$,
    $(\pi P)_u = \sum_{v \sim u} \frac{\deg(v)}{2m} \cdot \frac{1}{\deg(v)} = \sum_{v \sim u} \frac{1}{2m}
        = \frac{\deg(u)}{2m} = \pi_u$,
    since $u$ has $\deg(u)$ neighbours.
\end{proof}

\section{Convergence to the Stationary Distribution}

Convergence of vectors is entry by entry. We state the theorem without proof.

\begin{theorem}[Convergence of Random Walks]\label{thm:convergence-0910}
    Let $G$ be a connected, non-bipartite graph with $m$ edges. Then for every starting distribution
    $p_0$,
    \[p_t = p_0 P^t \longrightarrow \pi \qquad \text{as } t \to \infty, \qquad \text{where } \pi_v = \frac{\deg(v)}{2m}.\]
    So $\pi$ is the only stationary distribution of the walk.
\end{theorem}
